\section{Sharpness: a non-triangle ambiguity in \texorpdfstring{$\TCw\setminus\TCs$}{W-TC minus S-TC}}\label{sec:sharpness}

We now prove \cref{thm:sharpness} without appealing to any result outside this paper and its exact supplement.  The construction is a pair of four-leaf level-2 theta networks whose pendant placements differ.  Their standard semi-directed topologies are weakly tree-child but not strongly tree-child, and their JC images meet along a regular eight-dimensional open set.

\subsection{The four-leaf pair}

Let the internal vertices be a root $\rho$, tree vertices $A,B,D,E$, and reticulations $C,F$.  Both rooted networks have internal arcs
\begin{equation}\label{eq:theta-common-arcs}
\rho\to A,\quad \rho\to C,\quad A\to B,\quad B\to C,
\quad C\to D,\quad D\to E,\quad A\to F,\quad E\to F.
\end{equation}
Their pendant arcs are
\begin{align}
N:&\qquad B\to1,\quad D\to2,\quad F\to3,\quad E\to4,\label{eq:theta-N-pendants}\\
N':&\qquad E\to1,\quad D\to2,\quad F\to3,\quad B\to4.\label{eq:theta-Np-pendants}
\end{align}
After suppressing the root artifact, the nontrivial blob is a theta graph with three $A$--$C$ paths
\[
A-C,\qquad A-B-C,\qquad A-F-E-D-C.
\]
Its three simple cycles have lengths $3,5,6$, so both networks are binary and level~2, and each blob contains exactly one triangle.  The two standard reductions are shown in \cref{fig:theta-sharpness}.

\begin{figure}[p]
\centering
\begin{tikzpicture}[scale=.70,every node/.style={font=\scriptsize}]
\newcommand{\thetablob}[2]{%
\begin{scope}[shift={(#1,0)}]
\node[treev] (A#2) at (0,2.2) {$A$}; \node[treev] (B#2) at (-1.4,1.2) {$B$};
\node[reticv] (C#2) at (0,0) {$C$}; \node[treev] (D#2) at (1.3,-.8) {$D$};
\node[treev] (E#2) at (2.35,.1) {$E$}; \node[reticv] (F#2) at (2.0,1.55) {$F$};
\draw[reticedge] (A#2)--(C#2); \draw[ordinary] (A#2)--(B#2); \draw[reticedge] (B#2)--(C#2);
\draw[ordinary] (C#2)--(D#2); \draw[ordinary] (D#2)--(E#2); \draw[reticedge] (A#2)--(F#2); \draw[reticedge] (E#2)--(F#2);
\end{scope}}
\thetablob{-5}{L}
\thetablob{3.5}{R}
% left leaves
\node[leaf] (l1) at (-7.1,1.2) {$1$}; \node[leaf] (l2) at (-3.7,-1.8) {$2$}; \node[leaf] (l3) at (-3,2.35) {$3$}; \node[leaf] (l4) at (-1.85,-.55) {$4$};
\draw (BL)--(l1); \draw (DL)--(l2); \draw (FL)--(l3); \draw (EL)--(l4);
\node[font=\bfseries] at (-3.8,3.5) {$N$};
% right leaves swapped 1,4
\node[leaf] (r4) at (1.4,1.2) {$4$}; \node[leaf] (r2) at (4.8,-1.8) {$2$}; \node[leaf] (r3) at (5.5,2.35) {$3$}; \node[leaf] (r1) at (6.65,-.55) {$1$};
\draw (BR)--(r4); \draw (DR)--(r2); \draw (FR)--(r3); \draw (ER)--(r1);
\node[font=\bfseries] at (4.7,3.5) {$N'$};
\draw[<->,very thick] (-.75,1.2)--(.65,1.2) node[midway,above] {$\Theta$};
\node[align=center,text width=11.5cm] at (0,-3.2) {The pendant components at $B$ and $E$ are exchanged. The two standard semi-directed graphs are nonisomorphic and not $\Tmove$-equivalent. Each has five admissible rootings, exactly two tree-child, so both lie in $\TCw\setminus\TCs$.};

% leaf substitution schematic
\node[leaf] (p) at (-1.3,-5.2) {$i$};
\draw[-{Stealth[length=2.5mm]},thick] (-.5,-5.2)--(.8,-5.2) node[midway,above]{substitute};
\node[treev] (t) at (2,-5.2) {};
\node[leaf] (q1) at (3.25,-4.55) {$i_1$}; \node[leaf] (q2) at (3.25,-5.85) {$i_2$};
\draw (t)--(q1) node[midway,above] {$u$}; \draw (t)--(q2) node[midway,below] {$v$};
\node[align=center,text width=7cm] at (0,-7.15) {$\widetilde P(g_X,g_1,g_2)=P(g_X,g_1\oplus g_2)u^{\mathbf1[g_1\ne0]}v^{\mathbf1[g_2\ne0]}$};
\end{tikzpicture}
\caption{The Theta sharpness pair and identical leaf substitution. Repeated substitution gives a non-$\Tmove$ ambiguity in $\TCw\setminus\TCs$ for every $n\ge4$.}
\label{fig:theta-sharpness}
\end{figure}


\begin{proposition}[Tree-child class and topological separation]\label{prop:theta-class}
The supplied rooted representatives of $N$ and $N'$ lie in $\TCr$.  Each standard semi-directed topology has five admissible rootings, exactly two of which are tree-child.  Hence
\[
N,N'\in\TCw\setminus\TCs.
\]
The two standard semi-directed topologies are nonisomorphic and are not related by ordinary triangle redirection.
\end{proposition}

\begin{proof}
In the rooted networks \eqref{eq:theta-common-arcs}--\eqref{eq:theta-Np-pendants}, every internal vertex has a tree-or-leaf child and neither reticulation has a reticulation child.  Thus both chosen rootings are tree-child.

For the standard reductions, retain the arrowheads entering $C$ and $F$, suppress $\rho$, and enumerate each edge on which a root can be inserted without changing the retained arrowheads.  Propagating the forced ordinary orientations from the proposed root gives exactly five admissible rootings.  The root sites $A-C$ and $A-F$ give tree-child rootings; the remaining three give a tail of a retained arrow two reticulation children.  The same count and the same two root sites occur for both pendant placements.  This proves membership in $\TCw\setminus\TCs$.  The enumeration is reproduced from the arc lists by the independent checker in the supplement.

The unique triangle is $A-B-C$.  In $N$, leaf $1$ is adjacent to $B$, a triangle vertex; in $N'$, leaf $1$ is adjacent to $E$, which is not on a triangle.  This labelled graph property is preserved by isomorphism and by a redirection confined to the unique triangle.  Therefore the two topologies are neither isomorphic nor $\Tmove$-equivalent.
\end{proof}

\subsection{Exact common JC point}

Write every edge multiplier as $x_e\in(0,1)$ and let $\lambda_C,\lambda_F\in(0,1)$ be the two inheritance probabilities.  It is convenient to suppress the root and use the effective product on the root-split $A-C$ edge.  A rooted realization is recovered by the strict factorization
\begin{equation}\label{eq:theta-root-factor}
 x_{\rho A}=\frac23,\qquad x_{\rho C}=\frac34,
 \qquad x_{\rho A}x_{\rho C}=\frac12.
\end{equation}

For $N$, take
\begin{equation}\label{eq:theta-source-point}
\begin{gathered}
 x_{D2}=\frac12,\quad x_{DE}=\frac25,\quad x_{E4}=\frac38,
 \quad x_{EF}=\frac13,\quad x_{F3}=\frac12,\\
 x_{CD}=\frac9{20},\quad x_{AB}=\frac35,\quad x_{B1}=\frac15,\\
 x_{AC}=x_{AF}=x_{BC}=\frac12,
 \qquad \lambda_C=\lambda_F=\frac12.
\end{gathered}
\end{equation}
Let $\beta$ be the smaller real root of
\begin{equation}\label{eq:beta-polynomial}
43337075\beta^2-36083110\beta+7336259=0,
\qquad
\frac{441}{1250}<\beta<\frac{3529}{10000}.
\end{equation}
For $N'$, take
\begin{equation}\label{eq:theta-target-point}
\begin{gathered}
 x_{D2}=\frac12,\qquad x_{DE}=\frac{171}{775},\qquad
 x_{AF}=\frac{10339}{53010\beta},\\
 x_{AB}=\frac{24835\beta}{20678-24835\beta},\qquad
 x_{F3}=\frac{1767}{4832},\qquad x_{CD}=\frac{9934}{12215},\\
 x_{E1}=\frac{31}{190},\qquad x_{B4}=\frac{3}{20\beta},\\
 x_{AC}=x_{EF}=x_{BC}=\frac12,
 \qquad \lambda_C=\lambda_F=\frac12.
\end{gathered}
\end{equation}
The isolating interval in \eqref{eq:beta-polynomial} proves directly that every quantity in \eqref{eq:theta-target-point} lies strictly between zero and one.

The four-leaf JC orbit space has fourteen nonconstant coordinates.  In the ordering
\[
(A,B,C,D,E,F,G,H,J,K,L,M,N,O)
\]
used in \cref{app:theta-data}, the common point is
\begin{align}
(A,\ldots,H)&=
\left(
\frac{277}{8000},\frac3{40},\frac{67}{2400},\frac{127}{16000},
\frac{27}{5000},\frac{81}{5000},\frac{153}{160000},\frac9{500}
\right),\label{eq:theta-common-first}\\
(J,\ldots,O)&=
\left(
\frac{27}{16000},\frac{261}{320000},\frac{27}{10000},
\frac{27}{20000},\frac{153}{320000},\frac{183}{80000}
\right).\label{eq:theta-common-second}
\end{align}

\begin{proposition}[Exact Fourier equality]\label{prop:theta-equality}
Substitution of \eqref{eq:theta-source-point} and \eqref{eq:theta-target-point} into the displayed-tree formula \eqref{eq:fourier-network} gives identical values for all $256$ four-leaf Fourier coordinates, and hence for all $256$ pattern probabilities.
\end{proposition}

\begin{proof}
Coordinates with nonzero total character sum vanish on both sides.  The remaining coordinates reduce, under JC character symmetry, to the fourteen entries in \eqref{eq:theta-common-first}--\eqref{eq:theta-common-second}.  For the source, direct rational substitution gives those values.  For the target, clear the denominators and reduce every numerator difference modulo the quadratic polynomial in \eqref{eq:beta-polynomial}; each remainder is zero.  The complete list of orbit representatives and the exact sparse reductions appear in \cref{app:theta-data}, and the independent verifier evaluates the unreduced $256$-coordinate tensors directly.  The inverse Fourier transform is linear and invertible, so the pattern distributions agree as well.
\end{proof}

\subsection{A common regular eight-dimensional germ}

The equality at one point is not enough for the sharpness theorem.  We now show that the intersection contains a full-dimensional regular neighborhood.

Both parameterizations satisfy the six equations
\begin{align}
J-K-M+N&=0,\label{eq:theta-inv1}\\
J-AH-BF+CE&=0,\label{eq:theta-inv2}\\
GL-EN&=0,\label{eq:theta-inv3}\\
L^2-BEH&=0,\label{eq:theta-inv4}\\
BM-DL-B^2F+BCE&=0,\label{eq:theta-inv5}\\
BEO-BGH-CEL+DEH&=0.\label{eq:theta-inv6}
\end{align}
On the physical locus $BE>0$, five of these reconstruct $J,K,M,N,O$ from
$(A,B,C,D,E,F,G,H,L)$, while \eqref{eq:theta-inv4} gives
\[
L=+\sqrt{BEH}.
\]
Thus the physical main locus is a smooth sheet of the irreducible variety with localized coordinate ring
\begin{equation}\label{eq:theta-coordinate-ring}
\mathbb C[A,B,C,D,E,F,G,H,L,(BE)^{-1}]/(L^2-BEH),
\end{equation}
which has dimension eight.  Smoothness on the open locus follows from
$\partial(L^2-BEH)/\partial H=-BE\ne0$.

For the source, fix
\[
x_{AC}=x_{AF}=x_{BC}=\lambda_C=\lambda_F=\frac12
\]
and use the eight variables
\[
(P,s,Q,t,R,u,v,S)
=(x_{D2},x_{DE},x_{E4},x_{EF},x_{F3},x_{CD},x_{AB},x_{B1}).
\]
The Jacobian determinant of the map to $(A,\ldots,H)$ is
\begin{equation}\label{eq:theta-source-jac}
-\frac{P^3s^3Q^4tR^4u^3vS^3(s-1)^2(v-1)(v+1)^2}{16384},
\end{equation}
which is nonzero at \eqref{eq:theta-source-point}.  For the target, fix
\[
x_{AC}=x_{EF}=x_{BC}=\lambda_C=\lambda_F=\frac12
\]
and use
\[
(P',x,y,z,R',w,S',Q')
=(x_{D2},x_{DE},x_{AF},x_{AB},x_{F3},x_{CD},x_{E1},x_{B4}).
\]
The corresponding determinant is
\begin{equation}\label{eq:theta-target-jac}
-\frac{(P')^3x^2y^2z(R')^4w^4(S')^3(Q')^4
(x-1)^2(z-1)(z+1)^3}{32768},
\end{equation}
which is nonzero throughout the open target point \eqref{eq:theta-target-point}.

\begin{corollary}[Four-leaf regular overlap]\label{cor:theta-four-leaf}
The two four-leaf JC images share a relatively open regular eight-dimensional subset.  In particular, $N\bowtieJC N'$.
\end{corollary}

\begin{proof}
The six identities place both images in the same irreducible eight-dimensional locus \eqref{eq:theta-coordinate-ring}.  Equations \eqref{eq:theta-source-jac} and \eqref{eq:theta-target-jac} show that each image is locally an open eight-dimensional submanifold of the positive smooth sheet at their common point.  Their two image neighborhoods therefore have a nonempty relative intersection in that sheet.
\end{proof}

\subsection{The extra parameter defect}

The eight-dimensional result has one more generic fiber dimension than the standard local edge gauges alone predict.  This point matters for the all-$n$ dimension formula.

A rooted binary network with $n$ leaves and $r$ reticulations has
\begin{equation}\label{eq:sharp-edge-count}
E=2n+3r-2
\end{equation}
network edges; see \cref{app:dimensions}.  For $r=2$, the raw JC parameter count is
$E+r=2n+6$.  The ordinary root-split gauge and two local gauges at each reticulation account for five generic fiber dimensions, suggesting the upper value $2n+1$.  For the Theta topology there is one further generic dependence.  At $n=4$, the exact implicit description \eqref{eq:theta-coordinate-ring} and the nonzero determinants above prove dimension eight rather than nine.  Equivalently, the generic fiber has dimension six rather than five.  This extra one-dimensional defect persists under the leaf substitutions below because each substitution contributes two new identifiable arm multipliers and does not interact with the original fiber.

\subsection{Leaf substitution and the all-\texorpdfstring{$n$}{n} family}

Let $P$ be the Fourier tensor of either member of the four-leaf pair, written with one selected leaf port $a$.  Replace $a$ by a cherry with new leaves $1,2$ and edge multipliers $u,v\in(0,1)$.  If $g_X$ denotes the characters on all other leaves, then
\begin{equation}\label{eq:theta-leaf-substitution}
\widetilde P(g_X,g_1,g_2)
=P(g_X,g_1\oplus g_2)
 u^{\mathbf1[g_1\ne0]}v^{\mathbf1[g_2\ne0]}.
\end{equation}
For any nonzero $h\in G$, positivity gives
\begin{align}
uv&=\widetilde P(0,h,h),\label{eq:theta-cherry-product}\\
\frac uv&=
\frac{\widetilde P(g_X,h,0)}{\widetilde P(g_X,0,h)}
\quad\text{whenever }\bigoplus g_X=h.
\label{eq:theta-cherry-ratio}
\end{align}
Hence $u$ and $v$ are recovered uniquely by the positive square root, after which $P$ is recovered from \eqref{eq:theta-leaf-substitution}.  Leaf substitution is therefore a positive analytic embedding with a positive analytic inverse onto its image.  Replacing corresponding leaves of $N$ and $N'$ by the same rooted binary tree preserves the common tensor equality and adds exactly two model dimensions for every new leaf.

The standard semi-directed graphs remain nonisomorphic and non-$\Tmove$-equivalent because restricting to one chosen descendant leaf in every substituted component recovers the original labelled pair.  A tree-child rooting extends through every substituted rooted tree, while each original non-tree-child admissible rooting remains non-tree-child after substitution.  Thus membership in $\TCw\setminus\TCs$ is preserved.

Starting from the eight-dimensional four-leaf pair and applying $n-4$ cherry substitutions yields
\[
\dim \M_{N_n}^{\JC}
=\dim \M_{N_n'}^{\JC}
=8+2(n-4)=2n.
\]
The analytic inverse shows that the common open germ also gains two dimensions at every step.  This proves \cref{thm:sharpness} for every $n\ge4$.
